\documentclass{beamer}
\usepackage{color}
\usepackage{bm}
\usepackage{graphicx}
\usepackage{hyperref}
\usepackage{listings}
\beamertemplatenavigationsymbolsempty
\definecolor{mygreen}{rgb}{0,0.6,0}
\definecolor{mymauve}{rgb}{0.58,0,0.82}
\lstset{ %
  backgroundcolor=\color{white},   % choose the background color; you must add \usepackage{color} or \usepackage{xcolor}
  basicstyle=\footnotesize,        % the size of the fonts that are used for the code
  commentstyle=\color{mygreen},    % comment style
  deletekeywords={...},            % if you want to delete keywords from the given language
  extendedchars=true,              % lets you use non-ASCII characters; for 8-bits encodings only, does not work with UTF-8
  frame=single,                    % adds a frame around the code
  keywordstyle=\color{blue},       % keyword style
  language=Python,                 % the language of the code
  rulecolor=\color{black},         % if not set, the frame-color may be changed on line-breaks within not-black text (e.g. comments (green here))
  showspaces=false,                % show spaces everywhere adding particular underscores; it overrides 'showstringspaces'
  showstringspaces=false,          % don't show spaces in strings
  stringstyle=\color{mymauve},     % string literal style
  title=\lstname                   % show the filename of files included with \lstinputlisting; also try caption instead of title
}
\hypersetup{
    colorlinks=true,
    urlcolor=blue
}

\graphicspath{ {./img/} {./charts/} }


\title{Django's System Check Framework}
\author{Adam Johnson - me@adamj.eu}
\date{13 July 2015}

\begin{document}


\maketitle


\begin{frame}[fragile]\frametitle{What is it?}

    \begin{center}
        \includegraphics[width=10cm]{checkit}
    \end{center}

    \begin{itemize}
        \item It checks your code!
    \end{itemize}

\end{frame}


\begin{frame}[fragile]\frametitle{What is it?}

    \begin{itemize}
        \item Added to Django in version 1.7
        \item Automatically checks things after 'ready' time, such as models, deprecated features, and contrib apps
        \item Can be run alone as manage.py command
        \item Looks like this:
    \end{itemize}

        \begin{lstlisting}
$ ./manage.py check
System check identified some issues:

WARNINGS:

?: (1_6.W001) Some project unittests may not execute as
  expected.
    HINT: Django 1.6 introduced a new default test
    runner. It looks like this project was generated
    using Django 1.5 or earlier. You should ensure your
    tests are all running & behaving as expected. See https://docs.djangoproject.com/en/dev/releases/1.6/#discovery-of-tests-in-any-test-module
    for more information.
    \end{lstlisting}

\end{frame}


\begin{frame}[fragile]\frametitle{Why is it good?}

    \begin{center}
        \includegraphics[width=10cm]{celebrate}
    \end{center}

    \begin{itemize}
        \item Code quality!
    \end{itemize}

\end{frame}


\begin{frame}[fragile]\frametitle{Why is it good?}

    \begin{itemize}
        \item ``Beyond linting'' - code quality checks that can only be done at runtime, not by flake8
        \item Extensible - third party apps can provide checks too
        \item Runs on nearly every manage command - you can't forget to run it, unlike tests!
    \end{itemize}

\end{frame}


\begin{frame}[fragile]\frametitle{Using it}

    \begin{center}
        \includegraphics[width=10cm]{demo}
    \end{center}

    \begin{itemize}
        \item How do you use it?
    \end{itemize}

\end{frame}


\begin{frame}[fragile]\frametitle{Using it}

    \begin{itemize}
        \item Normally it does its work without you realizing, until you see an error and correct it
        \item However it's easy to write your own checks - the framework is very simple
        \item At YPlan, we already had some sanity checks in place as unit tests - converting them to checks has improved the situation
    \end{itemize}

\end{frame}


\begin{frame}[fragile]\frametitle{An ex-unit test}

    \begin{itemize}
        \item An example old sanity-checking unit test from the YPlan codebase:
    \end{itemize}

    \begin{lstlisting}
class SupportedFeaturesTest(SimpleTestCase):
    """Check features that might not be supported in
       some OS versions"""
    def test_strftime_dash_strips_leading_zeroes(self):
        # Not every strftime implementation has %- to
        # strip leading zeroes, but we use it
        dt = datetime(2000, 10, 11, 2, 12)
        self.assertEqual(
            dt.strftime('%-I:%M %P'),
            '2:12 am'
        )
    \end{lstlisting}

    \begin{itemize}
        \item Bad! The test might not be run; or if it is run and fails, its output is buried amongst all the other test failures of features that depend upon strftime's \texttt{\%-I}
    \end{itemize}
\end{frame}



\begin{frame}[fragile]\frametitle{Make it better!}

    \begin{center}
        \includegraphics[width=10cm]{better}
    \end{center}

    \begin{itemize}
        \item Let's make it better!
    \end{itemize}

\end{frame}


\begin{frame}[fragile]\frametitle{Writing a check}

    \begin{itemize}
        \item Easy - just \texttt{register} a function that returns a list of \texttt{Error} objects representing bad things:
    \end{itemize}

    \begin{lstlisting}
from django.core.checks import Error, Tags, register

@register(Tags.compatibility)
def check_system(app_configs, **kwargs):
    # app_configs is a list of AppConfig objects
    # kwargs is for future extension
    errors = []
    if bad_thing():
        errors.append(
            Error(
                "Stop! Bad thing has happened!",
                id='mycode.001'
            )
        )
    return errors
    \end{lstlisting}

    \begin{itemize}
        \item Put in \texttt{myapp/apps.py} next to your \texttt{MyAppConfig} class
    \end{itemize}

\end{frame}


\begin{frame}[fragile]\frametitle{Converting our unit test}

    \begin{itemize}
        \item Magicked into a sub-function:
    \end{itemize}

    \begin{lstlisting}
@register(Tags.compatibility)
def check_system(app_configs, **kwargs):
    errors = []
    errors.extend(_check_strftime())
    # ...
    return errors

def _check_strftime():
    errors = []
    dt = datetime(2000, 10, 11, 2, 12)
    if dt.strftime('%-I:%M %P') != '2:12 am':
        errors.append(Error(
            "strftime does not appear to support using "
            "%- to strip leading zeroes",
            id='yplan.E007'
        ))
    return errors
    \end{lstlisting}

\end{frame}


\begin{frame}[fragile]\frametitle{It's better!}

    \begin{itemize}
        \item Now the \texttt{strftime} check runs at \texttt{runserver} and most other \texttt{manage.py} commands
        \item If it fails, it's unignorable
        \item It takes less than a millisecond so doesn't really impact startup time
    \end{itemize}

\end{frame}


\begin{frame}[fragile]\frametitle{requirements.txt}

    \begin{center}
        \includegraphics[width=8cm]{phone-call}
    \end{center}

    \begin{itemize}
        \item Dev: "Hey Adam, code is failing with ImportError"
        \item Me: "...you didn't pip install requirements.txt like I said on skype ten minutes ago, did you? Does anyone listen to me??"
    \end{itemize}

\end{frame}


\begin{frame}[fragile]\frametitle{requirements.txt}

    \begin{itemize}
        \item A communication bottleneck, especially with nearly 100 dependencies!
        \item Worse when developing several branches concurrently, only one of which is based on the new version of a package
    \end{itemize}

    \begin{center}
        \includegraphics[width=8cm]{frozen}
    \end{center}

\end{frame}


\begin{frame}[fragile]\frametitle{Solution: a system check!!}

    \begin{itemize}
        \item At launch, compare \texttt{pip freeze} to parsed lines of \texttt{requirements.txt}
        \item Basically:
    \end{itemize}

    \begin{lstlisting}
from pip.operations.freeze import freeze as pip_freeze

def check_requirements(app_configs, **kwargs):

    with open(REQUIREMENTS_TXT_PATH, 'r') as f:
        requirements = pip_parse(f.read().split('\n'))

    current = pip_parse(pip_freeze())

    errors = []
    for name, version in requirements.iteritems():
        # Compare with current - if missing or wrong
        # version, append an Error() to errors...

    return errors
    \end{lstlisting}

\end{frame}


\begin{frame}[fragile]\frametitle{Solution: a check}

    \begin{itemize}
        \item Example output:
    \end{itemize}

    \begin{lstlisting}
$ ./manage.py runserver
CommandError: System check identified some issues:

ERRORS:
?: (yplan.E002) Package patchy has version 1.1.0 in
   requirements.txt but you have version 1.0.0 installed
   - you should `pip install -r ../requirements.txt`

System check identified 1 issue (0 silenced).
    \end{lstlisting}

    \begin{itemize}
        \item Will open-source this. Haven't solved git-based packages yet
    \end{itemize}

\end{frame}


\begin{frame}[fragile]\frametitle{Running checks as part of tests}

    \begin{itemize}
        \item By default Django has doesn't run checks during \texttt{manage.py test}
        \item This is because tests mutate settings, e.g. prefixing database names with \texttt{test\_} - checks can't be run before this is done
        \item But \textit{lots} of our API development is done on branches using \textit{only} the test suite. We need the checks to run in this case!
    \end{itemize}


    \begin{center}
        \includegraphics[width=7cm]{pointing}
    \end{center}

\end{frame}


\begin{frame}[fragile]\frametitle{Running checks as part of tests}

    \begin{itemize}
        \item Turns out to not be so hard with a custom test runner class (\texttt{settings.TEST\_RUNNER}):
    \end{itemize}

    \begin{lstlisting}
from django.core.management import call_command
from django_nose import NoseTestSuiteRunner


class YPlanTestSuiteRunner(NoseTestSuiteRunner):

    def run_suite(self, *args, **kwargs):
        call_command('check')
        return super(YPlanTestSuiteRunner, self) \
               .run_suite(*args, **kwargs)
    \end{lstlisting}

\end{frame}


\begin{frame}[fragile]\frametitle{Idea - ``Django strict mode''}

    \begin{itemize}
        \item PEP8 enforces a stricter subset of python for more consistent code; a checks app could enforce a stricter subset of Django features for more consistent applications
        \item Turn some of the `best practice recommendations' in the docs into actual rules, for example `url names should start with \texttt{appname-}'
        \item Any ideas?
    \end{itemize}

\end{frame}


\begin{frame}[fragile]\frametitle{Thank you}

    \begin{center}
        \includegraphics[width=10cm]{clapping}
    \end{center}

    \begin{itemize}
        \item \url{me@adamj.eu} ; blog at \url{http://adamj.eu/tech/}
    \end{itemize}

\end{frame}


\end{document}
